% answers_short.tex —— 2005 年·上海卷·文科 · 简短答案速查
% 由 tools/make_short_answers.py 自 answers.tex 的 \gksolhead/\gkanswer 提取，
% 并与 2005/2005_数学真题_上海文_答案要点.md 逐题比对后生成，勿手改（改源 MD 后重跑）。

\begin{gknotice}
\textbf{使用说明}：本节只给最终答案便于快速核对，完整推导见解析版。答案为编写者独立推导，\textbf{非官方标准答案}。
\end{gknotice}

\gkpart{填空题}
\begin{itemize}[leftmargin=2.6em]
  \item[\textbf{1.}] $f^{-1}(x)=4^x-1$
  \item[\textbf{2.}] $x=0$
  \item[\textbf{3.}] 最大值为 $11$
  \item[\textbf{4.}] $x+2y=4$
  \item[\textbf{5.}] $T=\pi$
  \item[\textbf{6.}] $-\dfrac{11}{14}$
  \item[\textbf{7.}] $\dfrac{x^2}{80}+\dfrac{y^2}{20}=1$
  \item[\textbf{8.}] $\dfrac37$
  \item[\textbf{9.}] $x+2y=2$
  \item[\textbf{10.}] $AC=3$
  \item[\textbf{11.}] $k\in(1,3)$
  \item[\textbf{12.}] U/PENDING。
\end{itemize}

\gkpart{选择题}
\begin{center}\begin{tabular}{ccccc}
\toprule
题号 & 13 & 14 & 15 & 16 \\
\midrule
答案 & A & A & B & C \\
\bottomrule
\end{tabular}\end{center}

\gkpart{解答题}
\begin{itemize}[leftmargin=2.6em]
  \item[\textbf{17.}] $\arccos\dfrac2{\sqrt5}$
  \item[\textbf{18.}] $z=-\dfrac12\pm\dfrac{\sqrt3}{2}i$
  \item[\textbf{19.}] (1) $k=1,b=2$；(2) $-3$
  \item[\textbf{20.}] (1) U/PENDING；(2) $2009$ 年底
  \item[\textbf{21.}] (1) $y^2=4x$；(2) $N\left(\dfrac85,\dfrac45\right)$；(3) $m<1$ 时直线与圆相交，$m=1$ 时相切，$m>1$ 时相离
  \item[\textbf{22.}] (1) $h(x)=\begin{cases}x-2,&x<1,\\-2x^2+7x-6,&x\ge1.\end{cases}$；(2) 最大值为 $\dfrac18$；(3) 可取 $f(x)=\cos2x+\sin2x$，$\alpha=\dfrac\pi4$
\end{itemize}
